\section{Công thức chung của tư bản (T – H – T')}
\subsection{Sự vận động của tiền trong lưu thông hàng hóa giản đơn}
Trong kinh tế chính trị Mác – Lênin, việc phân tích công thức \textbf{H–T–H (Hàng – Tiền – Hàng)} là tiền đề then chốt để làm rõ sự khác biệt giữa sản xuất hàng hóa giản đơn và sự vận động của tư bản. Ở hình thức này, tiền tệ chỉ đóng vai trò là phương tiện môi giới cho quá trình trao đổi.

Đặc trưng nổi bật nhất của lưu thông hàng hóa giản đơn chính là mục đích hướng tới \textbf{giá trị sử dụng}. Đặc trưng của lưu thông hàng hóa giản đơn là mục đích hướng tới giá trị sử dụng. Hàng hóa sau khi trao đổi sẽ rời khỏi lưu thông và đi vào lĩnh vực tiêu dùng để thỏa mãn nhu cầu của con người. Nếu hàng hóa không được tiêu dùng thì quá trình lưu thông chưa hoàn thành (C. Mác và Ph. Ăngghen, Toàn tập, t.23, tr.191–192).

\subsection{Sự vận động của tư bản theo công thức T – H – T'}
Khác với lưu thông hàng hóa giản đơn, tiền trong nền sản xuất tư bản chủ nghĩa vận động theo một logic hoàn toàn khác. Tài liệu kinh tế chính trị Mác – Lênin chỉ rõ ràng tiền trong nền sản xuất tư bản chủ nghĩa vận động theo công thức T – H – T', được gọi là công thức chung của tư bản. Trong công thức này, T là lượng tiền ứng ra ban đầu, H là hàng hóa được mua vào với tư cách là tư liệu trung gian, và T' là lượng tiền thu về sau quá trình lưu thông, trong đó 
\[
    T' = T + t
\] 
Với $t$ là phần giá trị trội ra lớn hơn, tức là giá trị thặng dư, và điều kiện bắt buộc là $t$ > 0.

Sự vận động theo công thức T – H – T' không phải là một hiện tượng ngẫu nhiên hay đặc thù của một ngành kinh doanh nào, mà là hình thức vận động đặc trưng và phổ quát của mọi hình thái tư bản trong nền kinh tế tư bản chủ nghĩa. Dù là tư bản công nghiệp, tư bản thương nghiệp hay tư bản cho vay, tất cả đều có điểm xuất phát là tiền (T) và điểm kết thúc là một lượng tiền lớn hơn (T'). Tư bản công nghiệp bỏ tiền mua máy móc, nguyên liệu và thuê lao động để sản xuất hàng hóa rồi bán lấy tiền nhiều hơn. Tư bản thương nghiệp bỏ tiền mua hàng rồi bán lại với giá cao hơn. Tư bản cho vay bỏ tiền ra cho vay rồi thu về tiền cùng lãi.

Mục đích của lưu thông tư bản không phải là giá trị sử dụng mà là \textbf{sự lớn lên của giá trị}. Như Mác đã nhận định: \textbf{"Sự vận động của tư bản là không có giới hạn"} (Toàn tập, t.23, tr.234), vì mục đích của nó là giá trị trao đổi trừu tượng chứ không phải là sự thỏa mãn nhu cầu cụ thể.

\subsection{Ý nghĩa của từng yếu tố trong công thức và bản chất vận động của tư bản}
Phân tích từng yếu tố trong công thức T – H – T' cho phép làm sáng tỏ bản chất kinh tế của quá trình vận động tư bản. Yếu tố T (tiền ứng ra ban đầu) đóng vai trò là giá trị ban đầu mà nhà tư bản đầu tư vào quá trình sản xuất hoặc lưu thông. Đây là điều kiện tiên quyết để toàn bộ chu trình kinh doanh bắt đầu. Yếu tố H (hàng hóa) đóng vai trò là phương tiện trung gian – khác hoàn toàn với vai trò của tiền trong lưu thông hàng hóa giản đơn. Hàng hóa ở đây được mua vào không phải để tiêu dùng mà để tạo điều kiện cho quá trình sản xuất và bán lại. Và yếu tố T' (tiền thu về) là lượng giá trị lớn hơn T, bao gồm T ban đầu cộng với giá trị thặng dư.

Tài liệu kinh tế chính trị Mác – Lênin định nghĩa rõ ràng: số tiền trội ra lớn hơn được gọi là giá trị thặng dư; số tiền ứng ra ban đầu với mục đích thu được giá trị thặng dư trở thành tư bản. Định nghĩa này có ý nghĩa lý luận cực kỳ quan trọng vì nó xác định bản chất của tư bản không phải là một vật phẩm cụ thể (tiền, máy móc, nguyên liệu) mà là một quan hệ xã hội – quan hệ giữa giá trị và quá trình tăng thêm giá trị. Tiền chỉ biến thành tư bản khi được dùng để mang lại giá trị thặng dư. Ngược lại, tiền dùng để mua sắm tiêu dùng cá nhân, dù nhiều đến đâu, cũng không phải là tư bản.

Từ đây có thể rút ra nhận định quan trọng: mục đích cuối cùng của tư bản không phải là hàng hóa mà là sự gia tăng giá trị. Đây là đặc trưng phân biệt giữa người tư bản với người sản xuất hàng hóa giản đơn. Nếu người sản xuất nhỏ lẻ tham gia trao đổi để thỏa mãn nhu cầu thiết thực, thì nhà tư bản tham gia vào mọi hoạt động kinh tế với một mục tiêu duy nhất: tích lũy giá trị ngày càng nhiều hơn. Chính động lực này đã tạo nên tính năng động phi thường của chủ nghĩa tư bản, nhưng đồng thời cũng tạo ra những mâu thuẫn nội tại sâu sắc mà Marx đã phân tích.

Sự gia tăng giá trị không ngừng nghỉ chính là bản chất vận động của tư bản. Không có điểm dừng tự nhiên cho quá trình này – khác với nhu cầu tiêu dùng có thể được thỏa mãn, sự tích lũy giá trị trong chủ nghĩa tư bản luôn hướng đến một T' lớn hơn, rồi T'' lớn hơn nữa. Đây là cơ sở để hiểu tại sao nền kinh tế tư bản chủ nghĩa luôn có xu hướng mở rộng sản xuất, thâm nhập thị trường mới và tìm kiếm những phương thức khai thác lao động ngày càng hiệu quả hơn.

\section{Mâu thuẫn trong công thức tư bản}
\subsection{Vấn đề lý luận đặt ra từ công thức T – H – T'}
Sau khi xác lập được công thức chung của tư bản T – H – T', C. Mác đặt ra một câu hỏi lý luận trọng tâm: giá trị thặng dư xuất hiện từ đâu trong quá trình vận động này? Đây không phải là câu hỏi đơn giản về mặt kỹ thuật kinh tế, mà là vấn đề mang tính quyết định cho toàn bộ học thuyết về tư bản và bóc lột. Việc trả lời câu hỏi này đòi hỏi một phân tích sâu sắc về bản chất của quá trình trao đổi và sản xuất trong nền kinh tế tư bản chủ nghĩa.

Điểm xuất phát của sự phân tích là thị trường – không gian mà ở đó mọi hàng hóa, bao gồm cả hàng hóa sức lao động, đều được mua bán. Theo nguyên tắc ngang giá, tức là nguyên tắc cơ bản của nền kinh tế hàng hóa, mỗi hàng hóa được trao đổi theo đúng giá trị của nó. Nếu nguyên tắc này được tuân thủ nghiêm ngặt, thì trong quá trình trao đổi, không có giá trị mới nào được tạo ra – người bán bán hàng theo đúng giá trị, người mua mua theo đúng giá trị, và tổng lượng giá trị trong xã hội không thay đổi. Điều này dẫn đến một kết luận mâu thuẫn: nếu trao đổi luôn ngang giá, thì giá trị thặng dư không thể xuất hiện từ quá trình trao đổi.
\subsection{Giới hạn của giả định trao đổi không ngang giá}
Nhận thức được sự bế tắc của giả định trao đổi ngang giá, người ta có thể đặt câu hỏi ngược lại: liệu giá trị thặng dư có thể xuất hiện nếu trao đổi diễn ra không ngang giá, tức là nếu nhà tư bản mua hàng hóa thấp hơn giá trị hoặc bán cao hơn giá trị?

Nếu người mua hàng hóa để bán hàng hóa đó cao hơn giá trị thì chỉ được lợi xét về người bán, nhưng xét về người mua thì lại bị thiệt. Trong nền kinh tế thị trường, mỗi người đều đóng vai trò là người bán và đồng thời cũng là người mua. Do đó, nếu được lợi khi bán thì lại bị thiệt khi mua. Đây là một lập luận có tính hệ thống và toàn diện: khi xem xét toàn bộ nền kinh tế như một chỉnh thể, thì tổng lợi ích từ việc bán giá cao sẽ bằng đúng tổng thiệt hại từ việc mua giá cao. Điều này có nghĩa là trao đổi không ngang giá chỉ tạo ra sự phân phối lại giá trị đã có trong nền kinh tế, chứ không tạo ra giá trị mới.

Có khẳng định rằng lưu thông – tức là mua và bán thông thường – không tạo ra giá trị tăng thêm xét trên phạm vi xã hội. Nhận định này có ý nghĩa lý luận sâu sắc. Nó phủ nhận quan niệm của các nhà kinh tế học trọng thương khi cho rằng lợi nhuận xuất phát từ thương mại và trao đổi. Ngay cả khi có những giao dịch cá biệt mà một bên được lợi hơn bên kia, điều đó cũng không giải thích được sự giàu lên của toàn bộ giai cấp tư bản với tư cách là một giai cấp xã hội. Sự làm giàu của giai cấp tư bản không thể giải thích bằng cách nói rằng họ liên tục lừa dối hay mua rẻ bán đắt, vì điều đó chỉ dẫn đến sự phân phối lại giá trị trong nội bộ các tầng lớp có vốn.
\subsection{Mâu thuẫn lý luận và cách Mác đặt vấn đề}
Từ hai vế phân tích trên, C. Mác đã xác lập được một mâu thuẫn lý luận sâu sắc trong công thức chung của tư bản: một mặt, giá trị thặng dư không thể xuất hiện từ lưu thông nếu trao đổi ngang giá; mặt khác, giá trị thặng dư cũng không thể tồn tại nếu không có lưu thông, vì chính thông qua lưu thông mà nhà tư bản thu được T' nhiều hơn T ban đầu. Mác đã tóm lược mâu thuẫn này một cách kinh điển: \textbf{"Tư bản không thể xuất hiện từ lưu thông và cũng không thể xuất hiện ở bên ngoài lưu thông. Nó phải xuất hiện trong lưu thông và đồng thời không phải trong lưu thông"} (C. Mác và Ph. Ăngghen: Toàn tập, t.23, tr.275).

Đây là một mâu thuẫn logic xuất hiện khi người ta cố gắng giải thích nguồn gốc của giá trị thặng dư chỉ dựa vào quá trình trao đổi – tức là lưu thông. Các nhà kinh tế học cổ điển trước Mác, như Adam Smith và David Ricardo, đã nhận ra sự tồn tại của lợi nhuận và giá trị thặng dư, nhưng không giải thích được bản chất thực sự của chúng. Họ thường quy nguồn gốc lợi nhuận cho các yếu tố không rõ ràng như sự khéo léo trong kinh doanh, sự may mắn hay năng suất của tư bản.

Cách Marx đặt vấn đề với mâu thuẫn này thể hiện tư duy biện chứng đặc trưng của ông. Thay vì cố gắng giải quyết mâu thuẫn bằng cách loại bỏ một trong hai vế, Marx chấp nhận cả hai: giá trị thặng dư phải được hình thành trong quá trình sản xuất (không phải trong lưu thông), nhưng đồng thời cũng phải được thực hiện thông qua lưu thông (bán hàng hóa trên thị trường). Điều này dẫn đến câu hỏi quyết định: trong quá trình sản xuất, điều gì tạo ra giá trị thặng dư? Câu trả lời mà Marx tìm ra chính là hàng hóa sức lao động – loại hàng hóa đặc biệt mà nhà tư bản mua được trên thị trường lao động.

Bí mật mà Marx phát hiện ra chính là sự tồn tại của một loại hàng hóa đặc biệt trên thị trường: \textbf{Hàng hóa sức lao động}. Theo Marx, nhà tư bản đã may mắn tìm thấy trên thị trường một loại hàng hóa mà \textbf{"bản thân giá trị sử dụng của nó có cái đặc tính riêng biệt là làm nguồn gốc sinh ra giá trị"} (C. Mác và Ph. Ăngghen: Toàn tập, t.23, tr.277). Việc phát hiện ra hàng hóa sức lao động chính là bước ngoặt lý luận vĩ đại, cho phép giải quyết triệt để mâu thuẫn trong công thức chung và mở ra chìa khóa để hiểu rõ toàn bộ cơ chế vận động của chủ nghĩa tư bản.

\section{Hàng hóa sức lao động và sự hình thành giá trị thặng dư}
\subsection{Điều kiện lịch sử để sức lao động trở thành hàng hóa}
Trong nghiên cứu về nền sản xuất tư bản chủ nghĩa, Karl Marx đã đưa ra một phát hiện mang ý nghĩa cách mạng trong lĩnh vực Kinh tế chính trị học Mác – Lênin, đó là chỉ ra bản chất của hàng hóa sức lao động. \textbf{"Sức lao động hay năng lực lao động là toàn bộ những năng lực thể chất và tinh thần tồn tại trong cơ thể, trong một con người đang sống, và được người đó đem ra vận dụng mỗi khi sản xuất ra một giá trị sử dụng nào đó"} (C. Mác và Ph. Ăngghen: Toàn tập, t.23, tr.277). Đây là một định nghĩa mang tính hệ thống và toàn diện, đặt nền móng cho toàn bộ phân tích tiếp theo về quan hệ lao động trong chủ nghĩa tư bản.

Sức lao động không nghiễm nhiên là hàng hóa trong mọi thời đại. Để quá trình mua bán sức lao động diễn ra trên thị trường, cần hội đủ hai điều kiện lịch sử cụ thể:
\begin{itemize}
    \item \textbf{Thứ nhất, về mặt pháp lý}: Người lao động phải được \textbf{tự do về thân thể}. Họ phải có quyền sở hữu đối với sức lao động của chính mình và chỉ bán nó trong một thời hạn nhất định. Marx nhấn mạnh rằng người lao động phải đối xử với sức lao động của mình như một hàng hóa, tức là "người chủ sở hữu vật báu của mình" (C. Mác và Ph. Ăngghen: Toàn tập, t.23, tr.278). Điều này phân biệt rõ rệt công nhân làm thuê với nô lệ hay nông nô – những người vốn thuộc quyền sở hữu của chủ nô hoặc bị lệ thuộc vào địa chủ.
    \item \textbf{Thứ hai, về mặt kinh tế}: Người lao động không có tư liệu sản xuất. Họ bị tước đoạt mọi phương tiện cần thiết để tự mình thực hiện quá trình lao động tạo ra sản phẩm. Đây là điều kiện mang tính quyết định, buộc họ phải đem bán "vật báu" duy nhất còn lại là sức lao động để đổi lấy tư liệu sinh hoạt.
\end{itemize}

Nghịch lý lịch sử này rất đáng chú ý: chính nền tảng pháp lý của tự do cá nhân – quyền sở hữu bản thân và quyền tự do ký kết hợp đồng – lại tạo ra cơ sở để hình thành quan hệ làm thuê và bóc lột. Người nông nô bị gắn chặt với đất đai bởi các ràng buộc pháp lý và cưỡng bức ngoài kinh tế, còn người công nhân tự do bị gắn chặt với nhà tư bản bởi sự cưỡng bức kinh tế – họ phải bán sức lao động không phải vì bị ép buộc về mặt thân xác, mà vì không có lựa chọn nào khác nếu muốn tồn tại. Đây là một trong những phát hiện sâu sắc nhất của Marx về bản chất của tự do trong xã hội tư bản.

\subsection{Giá trị của hàng hóa sức lao động}
Giống như mọi hàng hóa khác, hàng hóa sức lao động có hai thuộc tính cơ bản: giá trị và giá trị sử dụng. Việc phân tích hai thuộc tính này là nền tảng để hiểu được cơ chế hình thành giá trị thặng dư.
\subsubsection{Giá trị của hàng hóa sức lao động}
Marx chỉ rõ rằng giá trị của hàng hóa sức lao động cũng do thời gian lao động xã hội cần thiết để sản xuất và tái sản xuất ra sức lao động quyết định. Vì sức lao động tồn tại như một năng lực sống của con người, nên \textbf{"giá trị sức lao động là giá trị của những tư liệu sinh hoạt cần thiết để duy trì người sở hữu sức lao động đó"} (C. Mác và Ph. Ăngghen: Toàn tập, t.23, tr.283).

Về mặt cấu trúc, giá trị sức lao động được cấu thành bởi ba bộ phận hợp thành:
\begin{itemize}
    \item \textbf{Giá trị các tư liệu sinh hoạt}: Để duy trì đời sống bản thân người công nhân (ăn, mặc, ở, đi lại...).
    \item \textbf{Chi phí đào tạo}: Bao gồm các phí tổn huấn luyện để người lao động có được trình độ kỹ năng chuyên môn tương ứng với yêu cầu của sản xuất.
    \item \textbf{Giá trị các tư liệu sinh hoạt cho gia đình người công nhân}: Nhằm duy trì và tái sản xuất liên tục lớp người làm thuê cho thế hệ kế tiếp.
\end{itemize}
Một đặc điểm đặc thù giúp phân biệt sức lao động với hàng hóa thông thường là \textbf{yếu tố tinh thần và lịch sử}. Marx nhận định: \textbf{"Ngược lại với các hàng hóa khác, việc xác định giá trị sức lao động bao gồm một yếu tố tinh thần và lịch sử"} (C. Mác và Ph. Ăngghen: Toàn tập, t.23, tr.284). Quy mô và cơ cấu các nhu cầu thiết yếu không cố định mà biến đổi theo điều kiện địa lý, trình độ văn minh và cuộc đấu tranh của giai cấp công nhân ở mỗi quốc gia trong từng giai đoạn lịch sử. Điều này có nghĩa là "mức sống tối thiểu" là một phạm trù động, phản ánh sự phát triển của lực lượng sản xuất xã hội.

\subsubsection{Giá trị sử dụng của hàng hóa sức lao động – Chìa khóa giải quyết mâu thuẫn}
Nếu giá trị của sức lao động là đặc điểm quyết định giá cả của nó trên thị trường lao động, thì giá trị sử dụng của sức lao động mới là yếu tố quyết định khả năng sinh lời của tư bản. Đây chính là điểm then chốt mà Marx đã phát hiện để giải quyết mâu thuẫn trong công thức chung của tư bản.

Khác với các hàng hóa thông thường sẽ hao mòn hoặc mất đi giá trị sau quá trình tiêu dùng, giá trị sử dụng của hàng hóa sức lao động có đặc tính đặc biệt: \textbf{nó là nguồn gốc sinh ra giá trị}. Marx chỉ rõ rằng, quá trình tiêu dùng sức lao động chính là quá trình sản xuất ra hàng hóa và tạo ra giá trị mới. Khi nhà tư bản sử dụng sức lao động, họ không chỉ bảo tồn được giá trị của tư bản bất biến mà còn thu về một lượng giá trị lớn hơn chi phí mua sức lao động ban đầu. Marx khẳng định: \textbf{"Giá trị sử dụng của sức lao động... không phải là cái giá trị mà nó có, mà là cái giá trị mà nó tạo ra"} (C. Mác và Ph. Ăngghen: Toàn tập, t.23, tr.287).

Sự chênh lệch giữa giá trị mới tạo ra và giá trị của bản thân hàng hóa sức lao động chính là \textbf{giá trị thặng dư} (m). Ví dụ, nếu nhà tư bản trả 15 USD để mua sức lao động trong một ngày, người công nhân không chỉ tái sản xuất ra đúng 15 USD đó để bù đắp tiền lương mà còn làm việc thêm một khoảng thời gian (lao động thặng dư) để tạo ra giá trị vượt trội. Chính phần giá trị này là mục tiêu duy nhất của nền sản xuất tư bản chủ nghĩa.

Như vậy, hàng hóa sức lao động là chìa khóa giải quyết mâu thuẫn trong công thức chung của tư bản vì:
\begin{itemize}
    \item \textbf{Trong lưu thông}: Nhà tư bản tuân thủ nguyên tắc trao đổi ngang giá (mua sức lao động đúng giá trị của nó).
    \item \textbf{Trong sản xuất}: Nhà tư bản thu được giá trị thặng dư vì giá trị sử dụng của sức lao động có khả năng tạo ra giá trị mới lớn hơn giá trị của chính nó.
\end{itemize}
Phát hiện này của Marx đã đánh đổ các quan niệm sai lầm của kinh tế học tầm thường khi cho rằng lợi nhuận sinh ra từ việc "mua rẻ bán đắt" hay từ năng suất máy móc thuần túy. Nó vạch trần bản chất bóc lột tiềm ẩn ngay cả khi các bên thực hiện giao dịch "công bằng" trên thị trường.

\subsection{Quá trình sản xuất giá trị thặng dư và phân tích cụ thể}
Để làm sáng tỏ cơ chế hình thành giá trị thặng dư, Karl Marx đã sử dụng ví dụ về quá trình sản xuất sợi. Đây là cách tiếp cận khoa học giúp cụ thể hóa các khái niệm trừu tượng thành những minh chứng định lượng có tính thuyết phục cao.

Trong ví dụ này, giả định để tiến hành sản xuất, nhà tư bản ứng ra các khoản chi phí cho một ngày làm việc 8 giờ của công nhân như sau: 50 USD để mua 50 kg bông, 3 USD cho hao mòn máy móc để kéo số bông đó thành sợi và 15 USD để mua hàng hóa sức lao động (đây là mức giá ngang với giá trị sức lao động đã được thỏa thuận). Như vậy, tổng số tiền nhà tư bản ứng ra ban đầu là 68 USD.

Quá trình sản xuất được chia thành hai giai đoạn:
\begin{itemize}
    \item \textbf{Trong giai đoạn đầu (4 giờ đầu)}: Bằng lao động cụ thể, người công nhân chuyển giá trị của 50 kg bông (50 USD) và hao mòn máy móc (3 USD) vào giá trị sản phẩm. Đồng thời, bằng lao động trừu tượng, người công nhân tạo ra một lượng giá trị mới là 15 USD. Tại thời điểm này, tổng giá trị sản phẩm thu được là 68 USD. Nếu quá trình sản xuất dừng lại ở đây, nhà tư bản chỉ thu hồi đủ số vốn đã bỏ ra, không có giá trị thặng dư và số tiền ứng ra chưa thể biến thành tư bản. Marx gọi đây là \textbf{thời gian lao động tất yếu} – khoảng thời gian mà người lao động tạo ra một lượng giá trị vừa đủ để bù đắp giá trị sức lao động của chính mình.
    \item \textbf{Trong giai đoạn tiếp theo (4 giờ sau)}: Vì nhà tư bản đã mua quyền sử dụng sức lao động cho cả ngày (8 giờ), công nhân phải tiếp tục làm việc theo hợp đồng. Trong 4 giờ này, nhà tư bản bỏ thêm 50 USD mua bông và 3 USD hao mòn máy móc. Người công nhân tiếp tục tạo ra thêm 15 USD giá trị mới bằng lao động trừu tượng của mình. Tuy nhiên, vì tiền lương đã được trả trọn gói cho cả ngày, phần giá trị mới 15 USD này hoàn toàn thuộc về nhà tư bản. Đây chính là \textbf{thời gian lao động thặng dư}.
\end{itemize}
Kết thúc ngày làm việc, tổng chi phí nhà tư bản bỏ ra là 121 USD (bao gồm 100 USD bông, 6 USD máy móc và 15 USD tiền lương). Trong khi đó, tổng giá trị sản phẩm thu về là 136 USD. Khoản chênh lệch 15 USD chính là giá trị thặng dư (m). Marx kết luận rằng: \textbf{"Giá trị thặng dư chẳng qua chỉ là hiệu số giữa giá trị sản phẩm và giá trị của các yếu tố tiêu dùng để tạo ra sản phẩm đó"} (C. Mác và Ph. Ăng-ghen: Toàn tập, t.23, tr.312).

Qua phân tích này, bản chất của \textbf{Tư bản bất biến (c)} và \textbf{Tư bản khả biến (v)} hiện lên rõ nét:
\begin{itemize}
    \item \textbf{Tư bản bất biến (c)}: Là bộ phận tư bản (bông, máy móc) mà giá trị của nó được bảo tồn và chuyển nguyên vẹn vào sản phẩm, không thay đổi về lượng trong quá trình sản xuất.
    \item \textbf{Tư bản khả biến (v)}: Là bộ phận tư bản dùng để mua sức lao động. Nó không chỉ tái sản xuất ra giá trị của chính nó (v) mà còn tạo ra giá trị thặng dư (m). Đây là bộ phận duy nhất có sự biến đổi về số lượng, xác lập công thức tổng quát của giá trị hàng hóa:
    \[
        G = c + v + m 
    \]
\end{itemize}
Như vậy, giá trị thặng dư thực chất là bộ phận giá trị mới dôi ra ngoài giá trị sức lao động, là kết quả của lao động không công mà người công nhân đã thực hiện cho nhà tư bản dưới sự che chắn của hợp đồng lao động tự nguyện.

\subsection{Ý nghĩa lý luận của quy luật giá trị thặng dư}
Quy luật giá trị thặng dư không chỉ là một học thuyết kinh tế đơn thuần mà còn là công cụ nhận thức mang tầm vóc triết học và xã hội học sâu sắc, giúp nhìn thấu bản chất của toàn bộ hệ thống tư bản chủ nghĩa. Ph. Ăng-ghen đã đánh giá phát hiện này của Marx là một trong hai phát kiến vĩ đại nhất (cùng với duy vật lịch sử) đã đưa chủ nghĩa xã hội từ không tưởng trở thành khoa học.

zÝ nghĩa của quy luật này được thể hiện qua ba tầng bậc chính:
\begin{itemize}
    \item \textbf{Tầng bậc 1: Giải thích khoa học nguồn gốc của lợi nhuận}. Trước Marx, các nhà kinh tế học cổ điển thường quy lợi nhuận cho sự tiết kiệm hay năng lực quản lý. Marx vạch rõ rằng lợi nhuận thực chất là hình thái biến tướng của giá trị thặng dư – phần giá trị do lao động của công nhân tạo ra nhưng bị nhà tư bản chiếm đoạt. Ông khẳng định: \textbf{"Sản xuất ra giá trị thặng dư hay lợi nhuận, đó là quy luật tuyệt đối của phương thức sản xuất này"} (C. Mác và Ph. Ăng-ghen: Toàn tập, t.23, tr.871).
    \item \textbf{Tầng bậc 2: Vạch trần bản chất bóc lột tiềm ẩn}. Điểm đặc sắc của lý luận này là chứng minh bóc lột vẫn tồn tại ngay cả khi trao đổi diễn ra hoàn toàn ngang giá. Khác với sự cưỡng bức thân xác trong chế độ phong kiến, bóc lột tư bản chủ nghĩa là sự cưỡng bức kinh tế thuần túy, được che đậy tinh vi dưới hình thức hợp đồng tự nguyện.
    \item \textbf{Tầng bậc 3: Nền tảng cho toàn bộ hệ thống lý luận Mác-Lênin}. Học thuyết này là cơ sở để phân tích sự tích lũy tư bản, khủng hoảng kinh tế và sự phân chia giai cấp. Nó cho thấy sự vận động của tư bản là một chu trình liên tục: \textbf{T – H (SLĐ, TLSX)… SX … H’ – T’}. Để thực hiện được giá trị thặng dư, hàng hóa phải được bán đi trên thị trường; nếu không, chu trình sẽ đứt gãy và dẫn đến khủng hoảng.
\end{itemize}
\textbf{Về mặt thực tiễn}, đối với người lao động, hiểu rõ quy luật này giúp họ nhận thức đúng đắn về địa vị của mình. Tiền công thực tế chỉ là một phần nhỏ giá trị họ tạo ra; phần còn lại chính là nguồn gốc sự giàu có của chủ sở hữu. Nhận thức này là tiền đề để người lao động bảo vệ lợi ích chính đáng và đấu tranh cho sự công bằng trong quan hệ lao động.